\chapter{The Sound, and What It Took}\label{cha:sound}

Every other chapter of this book describes the machine as it is. This
one describes how it got that way, because the sound took longer to get
right than anything else in it and the route was instructive. Nothing
here is needed to use the machine. It is here because a fantasy computer
is a series of decisions, and the decisions about sound were the ones
that most often turned out to be about something else.

It has one thesis, and it is worth stating before the evidence rather
than after: \textbf{the sound was hard because keeping a real-time audio
ring fed from a machine that runs in bursts is hard, not because
emulating a SID is expensive.} Every fault the sound reported turned out
to be somewhere else, and the chip emulation --- the part everyone
suspects --- has never once been on the critical path.

If you are reading it to find out what happened to the SID chips: they
are gone, and the last two sections say why. The reason is not the one
you would guess from the six weeks described here, which is most of what
makes it worth writing down.

\section{Four chips, because it is a fantasy}
The C64 had one SID. The plan for this machine was four, at
\texttt{\$D400} and every \texttt{\$20} after it, for no better reason
than that a machine which never existed may have whatever its author
thinks it should have and four is a chord. Dag Lem's reSID --- the
emulation that has been the reference for thirty years --- was vendored
unchanged from VICE~3.3, wrapped in \pth{core/sid.cc}, and mixed to
16-bit mono. The frontend fed a ring buffer a scanline at a time and an
audio callback drained it at 48\,kHz.

That took an afternoon and it worked. Everything after this section is
the six weeks of it not quite working.

\section{A SID is not a CPU}
The first mistake was in the first version and is the kind only a
fantasy machine can make: the SIDs were clocked at the CPU's clock.

On a C64 the two are the same crystal, so the question never comes up.

Here the CPU runs at 40.5\,MHz and the SID is a 1\,MHz part, and
clocking reSID at 40.5 million cycles a second cost \textbf{21
milliseconds a frame} --- more than the whole frame --- and put every
pitch out by a factor of forty-one. The fix was one line and the lesson
outlived it: \emph{the SID's clock is not the CPU's}, and every later
change that touched the clock had to be told so again. The most recent
time was in August 2026, in a benchmark that swept the CPU clock and
forgot to tell reSID, and rendered every note at the wrong rate.

\section{``Choppy choppy choppy'' --- which was not the sound}
The machine's first outing on real hardware, a Raspberry Pi 3B+ on a
television, produced the report that named the problem for the rest of
the project: \emph{the SID sounds on the RPi3 are very very bad, choppy
choppy choppy}.

It was not an audio bug. The Pi was running the machine at about
17\,ms against a 16.7\,ms frame, and two recent changes had spent that
slack --- a debug recorder storing every instruction fetch, and a rule
that made the I/O page win on every memory access, which put an extra
comparison on every read and write in the machine. Neither was visible
on a desktop. Both were fatal on the Pi, and what they broke was the
sound, because the sound is the first thing that breaks when a frame
runs long.

That is the pattern the rest of this appendix is about. \textbf{The
sound is the machine's alarm bell.} Almost every time it rang, the fault
was somewhere else.

\section{Silence is not free}
With a clean power supply and honest measurements, the Pi at 20\,MHz sat
1.6\,ms over its frame budget. Of the 18.3\,ms, \textbf{3.4 was the
SIDs} --- all four of them, three of them silent.

reSID costs the same for silence as for sound: a chip that has been
written to is a chip that must be clocked, whether or not anything is
audible. So a chip is now clocked and mixed only once the machine has
written to it since reset, which at the prompt is three chips' worth of
frame given back. Note what that 3.4\,ms was --- four chips
\emph{clocked}, not four chips playing --- and that it is the largest
figure the sound has ever cost anywhere in this project.

That fix has a tail which is still in the machine: the cost is latched
by the \emph{first write}, not by whether anything is sounding. A
program that pokes a SID register during setup and then goes quiet pays
the full price for the rest of the session. On a desktop that is half a
millisecond of sixteen and nobody notices. On a Pi with no slack it is
worth knowing.

\section{The ring that started empty}
On the Pi the audio callback is served by the same core that runs the
machine, once a frame, a block at a time. A frame makes 800 samples; the
block was 1024; the ring started empty. So the ring's level lived
between nothing and one frame, and a callback regularly found 800 real
samples and 224 of silence --- a gap, on a steady beat, at a frame rate
the machine was comfortably holding.

Smaller blocks and a lead of three of them --- 32\,ms, which a music
player does not feel --- and the level never touches the floor.

\textbf{It made things worse.} At 15\,MHz the gap count went from 40--48
in two seconds to 52--62.

\section{Chips that play on without the machine}
The lead failed because the model behind it was wrong, and the numbers
said so plainly. At 15\,MHz the Pi ran at 58--60 frames a second, not a
locked 60. Audio was made only by the machine's frames, 800 samples
each. At 58 frames that is 46,400 samples a second against the 48,000
the device consumes. \emph{Any} shortfall drains the ring and then gaps
for ever after; a lead only postpones the first one. Even rounding the
cycles per scanline down from 520.83 to 520 is a 0.16\,\% deficit on its
own.

So the machine stopped tying its sound to its frames. Real hardware does
not stop its sound chips because the processor stalled --- a SID keeps
oscillating whatever the 6502 is doing --- and now nor does this one:
after a frame, if the ring is below its target, the SIDs are clocked on
\emph{without the CPU} until it is full. The SID clock does not move, so
pitch does not move; a slow frame sustains a note a fraction longer
rather than cutting it.

This is the most important decision in the whole story, and it is the
one that came back.

\section{The meter that read zero}
With the chips decoupled, and with the Pi's callback served every 128
lines rather than once a frame, the gap counter read \textbf{zero}, on
three boots, at the default clock. By every instrument the project had,
the sound was clean.

Doc: \emph{both say 0 drops but I still can hear them.}

He was right and the meter was wrong. The gap counter counts callbacks
that found the ring \emph{empty} --- it counts silence. And the
decoupling exists precisely to ensure the ring is never empty. So across
the entire band where a host is merely losing rather than drowning, the
counter reads zero and there is nothing wrong with the sound except that
a growing fraction of it was not made by the program. The SIDs were
playing on without the machine. That is what he was hearing.

It is now counted. Beside the gaps there is a second number: how many
samples the machine did not make and the chips supplied. On one host,
at rows both instruments had called clean:

\begin{center}
\small
\begin{tabular}{@{}l l l r@{}}
\toprule
\textbf{Clock} & \textbf{Frames} & \textbf{Gaps} & \textbf{Filled in} \\
\midrule
30\,MHz & 53 fps & 0 & 36{,}464 \\
20\,MHz & 60 fps & 0 & 16{,}818 \\
10\,MHz & 60 fps & 0 & 1{,}609 \\
\bottomrule
\end{tabular}
\end{center}

\noindent of 96{,}000 samples in the two seconds measured. At 30\,MHz
\emph{more than a third of the sound was not the program's}, and every
meter said clean.

\emph{Clean} now means three things at once, because any two of them can
be true while the sound is wrong: sixty frames a second, no gaps, and
under two per cent filled in.

\begin{aside}
\textbf{The moral, three times over.} A performance regression that
presented as an audio bug. An audio fix that made the gap count worse. A
gap counter that read zero because an earlier fix had guaranteed it
would. Each time, the instrument was measuring something adjacent to the
fault --- and each time, the person listening to the machine was right
before the numbers were.
\end{aside}

\section{A machine keeping somebody else's time}
The newest finding, and the one that reframes the rest: on the desktop
the machine was not keeping its own time at all.

The Pi has always paced itself against the wall clock, because its
graphics layer blocks until the display flips and a frame that overran
by a millisecond would wait for the next one and halve the frame rate.
The desktop did the opposite --- it left the pacing to the display's
vertical sync and did none of its own, so the machine ran at whatever
the compositor handed it. On one host, a 60\,Hz display, that was
\textbf{51.8 frames a second}, with the act of putting the picture on
the glass taking 16.9\,ms of a 19.3\,ms frame.

A machine at 51.8 frames a second makes 51.8 frames of sound where the
device wants 60. The chips filled in the missing seventh, exactly as
designed, and it was audible --- and the new counter measured it at
17\,\%. The machine now keeps its own 60\,Hz on the desktop as it always
has on the Pi, and presents when it is ready:

\begin{center}
\small
\begin{tabular}{@{}l r r@{}}
\toprule
& \textbf{Whole frame} & \textbf{Onto the glass} \\
\midrule
before & 19.294\,ms (51.8 fps) & 16.867\,ms \\
after  & 16.668\,ms (60.0 fps) & 0.260\,ms \\
\bottomrule
\end{tabular}
\end{center}

\section{Where it stood, when it was decided}
By the end of it the machine had four reSID chips in it and ran as many
of them as you asked for, and the question of whether to keep them was
open. It had been raised as a performance question, after a Raspberry Pi
that would hold only 15\,MHz. On that ground it was settled, and settled
\emph{against} replacement. \textbf{reSID was not the cost.} One chip
sounding measured \textbf{0.039\,ms of a 16.67\,ms frame} on a desktop
--- a quarter of one per cent --- and four chips written since reset came
to about half a millisecond. The largest the sound ever cost was the
3.4\,ms of \S C.4, on the Pi, before silent chips stopped being clocked,
and that was a bookkeeping fault rather than a fidelity one. Nothing
measurable was to be bought by removing them.

Where the question was real was identity. The SID is a part this machine
inherited rather than chose: it was in here because the C64 had one, in
the same way the \verb|*| prefix is in here because Acorn had it. A
machine that never existed is entitled to a voice its author picked on
purpose --- and the original specification for this one read \emph{4
SIDs\slash OPL2}, so a synthesiser of the other tradition would not be a
betrayal of the plan but the rest of it. Against that: four SIDs were
what the machine had sounded like since its first afternoon, they cost
almost nothing, and the entire body of C64 music played on them, which
is not a small thing to hand back.

\section{What was decided}
On 1~September~2026 the SIDs were turned \emph{off} --- not deleted. The
OPL2 became the machine's default chip on every host, the menu rows that
offered SIDs went away, and one line in \texttt{k4510.cfg} gave them
back. That is the disposition this project calls OFF, and it exists so
that a decision of this size can be tried before it is taken.

On 5~September it was taken: \emph{``nuke anything having to do with
SIDs''}. Out went reSID, the four-chip mix with its phase carry, the SID
clock select, the settings and their migration, the SID demos and the
player, the tune archive, and every register write to \texttt{\$D400} in
the ROM, in EhBASIC and in Mad Pascal's runtime. The machine has one
sound chip now and it is the OPL2, at \texttt{\$D480}, nine FM voices
wired the AdLib's way.

Three things fell out of it that are worth recording:

\begin{itemize}[leftmargin=1.4em]
\item \textbf{The sound sequencer had to learn FM.} It was the SIDs' last
  real user --- BBC BASIC's \verb|SOUND|, Mad Pascal's \verb|Sound| and
  the invaders all reach it --- so its four channels became OPL2 voices,
  with the Beeb's noise channel as a feedback-heavy patch and the other
  three as plain two-operator tones. The patch is rewritten on every
  note, because a program may have zeroed the chip for its own purposes
  and the sequencer must still sound after it.
\item \textbf{The render seam got a name.} \pth{core/audio.c} is a page
  of C that owes nothing to any chip: cycles in, samples out at the
  device's rate, the clock advanced. Everything this appendix is about
  --- the ring, the lead, the filling-in --- lives there now, and a
  future chip plugs into it rather than into the machine.
\item \textbf{Save states changed format}, and old ones refuse to load
  rather than pretending.
\end{itemize}

\section{And the moral, once more}
Read back, the six weeks argue for a conclusion the last two sections
then decline to draw. Every fault the sound reported was somewhere else;
the chip emulation was never on the critical path; and when the machine
was measured properly, the SIDs cost a quarter of one per cent of a
frame. The performance case for removing them was nil, and this appendix
says so at length.

They were removed anyway, and rightly, because the question was never
performance. It was what the machine is. That is the difference between
a fantasy computer and a product: a product's parts are decided by
measurement, and a fantasy computer's are decided by its author ---
having first taken the trouble to find out what the measurement actually
said, so that the decision is an honest one rather than a rationalised
one.

Everything above this section is kept exactly as it was written, because
the route matters more than the destination, and a book that quietly
rewrote its own history would be worth less than one that admits what it
tried.
